\documentclass[11pt]{article}
\usepackage{icmlko}
\begin{document}

\title{한 칸, 도메인마다 다른 재료 --- 네 번 막힌 벽을 옆으로 지나간다}
\author{월드모델 실험실}
\date{노트 449 · 2026년 8월 2일}
\maketitle

\begin{abstract}
같은 벽을 네 번 만났다. 시장팝업 전용 축은 그 도메인을 $+0.0241$ 올리는데
판엔 못 온다(노트 412). 표지 둘째는 판 $+0.0035$ 인데 전이를 $-0.1140$
잃는다(노트 420). 분기는 애니를 $+0.0167$ 올리는데 판은 $+0.0003$
이다(노트 448). 진단은 늘 같았다 --- \textbf{열을 하나 더하면 그 열이
``이 행은 애니다''라고 말하는 데 절반을 쓴다.}

이번엔 열을 안 늘렸다. \textbf{한 칸을 두고 도메인마다 다른 재료를
넣었다.} 도메인마다 자기 \emph{전용} 축 중 학습 라벨과 제일 붙은 것을 골라
\textbf{도메인 안에서 순위 정규화}해 같은 칸에 넣는다(부호는 맞춰서).
뜻이 ``이 행은 자기 도메인의 특기 축에서 몇 등인가''라 \textbf{정보는
나르고 정체는 안 나른다}.

\begin{center}
\begin{tabular}{lrr}
\hline
 & 판 차 & 부호 \\
\hline
전 축에서 고름 & $-0.0164$ & --- \\
전용 $\le 2$ 도메인 & $+0.0025$ & $7/12$ \\
\textbf{전용 $\le 4$ 도메인} & $\mathbf{+0.0046}$ & $\mathbf{10/12}$ \\
\hline
\end{tabular}
\end{center}

\textbf{채택.} 날짜 통제 판도 $+0.0062 \cdot 11/12$ 로 같이 오르고(노트
446 가드), LODO 도 안 막는다(문턱 넘은 다섯 중 오름 넷 · 평균 $+0.0202$).
챔피언은 \textbf{판 $0.4811 \to 0.4856$}, 날짜 통제 판 $0.5195 \to
0.5256$.
\end{abstract}

\begin{figure}[t]
\centering
\includegraphics[width=\textwidth]{figs/f1.pdf}
\caption{\textbf{가} --- 두 번 틀리고 세 번째에 됐다. \textbf{나} ---
왜 두 번 틀렸나. 칸이 빈 도메인이 세금을 문다. 칸을 더 채우니 세금이
$-0.0079 \to -0.0024$ 로 준다.}
\end{figure}

\section{두 번 틀린 자리가 답을 줬다}

\textbf{첫 번째 --- 축 고르기를 전 축에 열었다.} 열하나 중 \textbf{일곱이
\texttt{target\_breadth}} 를 골랐다. 이미 모두가 쓰는 공유 축이다. 그러니까
장치가 판 것은 ``전용 정보를 나르는 칸''이 아니라 ``공유 축의 도메인 안
순위 사본''이었고, 그건 판을 $\mathbf{-0.0164}$ 떨어뜨린다. \emph{재려던
물음을 안 재고 있었다.}

\textbf{두 번째 --- 전용을 $\le 2$ 도메인으로 좁혔다.} 이번엔 옳은 것을
쟀는데 관문을 못 넘었다($+0.0025 \cdot 7/12$). 그런데 갈라 보니 답이
있었다.

\begin{center}
\begin{tabular}{lrr}
\hline
 & 칸 찬 도메인 & 칸 빈 도메인 \\
\hline
전용 $\le 2$ (6 채움) & $+0.0054$ & $-0.0079$ \\
전용 $\le 4$ (7 채움) & $\mathbf{+0.0147}$ & $-0.0024$ \\
\hline
\end{tabular}
\end{center}

\textbf{칸이 빈 도메인이 세금을 문다.} 빈 칸은 (값 $0.5$, 표시자 $0$)이
되는데 \textbf{나무가 그 표시자로 가른다} --- ``특기 축이 있는 무리 대
없는 무리''로. \emph{쫓아낸 도메인 이름표가 뒷문으로 다시 들어온 것이다.}

진단이 처방을 줬다. \textbf{더 많이 채우면 세금이 준다.} 문턱을 $\le 4$
로 풀어 하나 더 채우니 세금이 $-0.0079 \to -0.0024$ 로 줄고, 찬 도메인의
이득은 $+0.0054 \to +0.0147$ 로 는다. 관문을 넘는다.

\section{아이돌 하나가 끄는 게 아니다}

관문에서 아이돌이 $+0.0706$ 으로 제일 크지만 유보가 $51$행이라 판 무게가
$1.5\%$ 다. 무게로 곱하면 $+0.0011$ --- 판 이득 $+0.0046$ 의 넷의 하나다.
나머지는 \textbf{애니}($+0.0119 \cdot 12/12$ · 무게 $18\%$)와
\textbf{웹툰}($+0.0086 \cdot 10/12$ · 무게 $21\%$)이 낸다. \emph{무게가
큰 두 도메인이 부호도 제일 곧다.}

전이도 같은 쪽이다. LODO 에서 \textbf{애니가 $+0.0279 \to +0.2093$} 이다
--- 안 본 도메인으로 채점하는데도 남들이 특기 칸을 갖는 게 애니를 돕는다.

\section{구현}

\texttt{F18\_bagboost.SPEC}. \texttt{fit} 이 학습 라벨만 보고 도메인마다
축을 고르고(\texttt{\_spec\_pick}), \texttt{fit} 과 \texttt{predict} 이
같은 축으로 칸을 채운다(\texttt{\_spec\_col}). \textbf{안 배운 도메인은
고른 축이 없어 칸이 저절로 빈다} --- 집 밖 예보에서 알아서 꺼진다.

가드: 켜고 끈 판이 관문 수치와 맞고($+0.0045$ 대 $+0.0046$, 통제 $+0.0060$
대 $+0.0062$), 입력을 안 건드리고, 모르는 도메인 이름에는 칸이 빈다.

되돌림 조항(노트 432): \textbf{``특기 칸 없는 F18''}.

\section*{남은 것}

\begin{itemize}
\item \textbf{네 도메인이 아직 빈다}(도서 · 만화 · 시장팝업 · 팝업). 세금이
$-0.0024$ 로 줄었지 $0$ 이 아니다. 다 채우면 $0$ 이 되나 --- 문턱을 더
풀면 ``전용''이 아니게 된다. 다른 채움법이 필요하다.
\item \textbf{칸을 둘로.} 지금은 도메인마다 축 \emph{하나}만 나른다.
둘째 칸이 더 버나, 아니면 다시 이름표가 되나.
\item \textbf{집 밖에서도 채울 수 있나.} 노트 443 의 몇 샷과 같은
꼴이다 --- 새 수집의 라벨 마흔이면 그 수집의 특기 축도 고를 수 있다.
\item LODO 씨앗 셋$\to$아홉 --- 열 노트째 미룸.
\end{itemize}

\end{document}
